\documentclass[10pt,a4paper]{extarticle}
\usepackage[utf8]{inputenc}
\usepackage[T1]{fontenc}
\usepackage[spanish,es-noshorthands]{babel}
\usepackage[left=15mm,right=15mm,top=11mm,bottom=11mm]{geometry}
\usepackage{lmodern,microtype,amsmath,booktabs,tabularx,tikz,xcolor,hyperref}
\usetikzlibrary{arrows.meta,positioning}
\definecolor{azul}{HTML}{173B5E}
\definecolor{oro}{HTML}{C58A2A}
\definecolor{turquesa}{HTML}{087C84}
\definecolor{tinta}{HTML}{253645}
\definecolor{gris}{HTML}{EDF2F5}
\definecolor{rojo}{HTML}{B55A4D}
\hypersetup{colorlinks=true,urlcolor=turquesa,linkcolor=azul}
\pagestyle{empty}
\setlength{\parindent}{0pt}
\setlength{\parskip}{2.4pt}
\setlength{\abovedisplayskip}{5pt}
\setlength{\belowdisplayskip}{5pt}
\setlength{\abovedisplayshortskip}{4pt}
\setlength{\belowdisplayshortskip}{4pt}
\renewcommand{\arraystretch}{1.12}
\color{tinta}
\newcommand{\seccion}[2]{\vspace{3pt}{\color{azul}\bfseries\large #1\enspace #2}\par\vspace{1pt}{\color{oro}\rule{12mm}{1.2pt}}\par\vspace{1pt}}
\begin{document}
{\sffamily\fontsize{8.5}{10}\selectfont\bfseries\color{azul}
UNIVERSIDAD DE CONCEPCIÓN \enspace/\enspace FACULTAD DE INGENIERÍA \enspace/\enspace DEPARTAMENTO DE INGENIERÍA INDUSTRIAL}\par
\vspace{3pt}
{\color{oro}\rule{\textwidth}{1.4pt}}\par
\vspace{4pt}
{\color{azul}\fontsize{17}{19}\selectfont\bfseries Formulación automática de modelos de optimización}\par
{\color{azul}\large Adaptación QLoRA y evaluación frente a la línea base}\par
\vspace{4pt}
{\small\textbf{Entrega 2 | Inteligencia Artificial Generativa (580694)}}\par
{\small Rayen Muñoz, Isidora Rivera, Sebastián Soto y José Luis Erices}\par
\vspace{4pt}
\noindent\colorbox{gris}{\parbox{\dimexpr\textwidth-2\fboxsep\relax}{\small
\textbf{\color{azul}Por qué importa.} Omitir la población al localizar centros cambia a quién prioriza el objetivo; omitir el vínculo de apertura permite atender sectores desde centros cerrados. El sistema produce borradores paramétricos para revisión experta.}}
\vspace{3pt}

{\color{azul}\bfseries\large Pipeline del experimento}\par
\vspace{-1pt}
\begin{center}
\begin{tikzpicture}[>=Stealth,font=\fontsize{8.1}{9.2}\selectfont,
caja/.style={draw=azul,rounded corners=2pt,align=center,minimum width=37mm,minimum height=10mm,fill=white,inner sep=2pt},
dest/.style={caja,fill=gris},
flecha/.style={-Stealth,draw=azul,line width=.75pt}]
\node[dest] (datos) at (1.9,0) {\textbf{1050 pares}\\enunciado + formulación};
\node[dest] (ajuste) at (6.1,0) {\textbf{QLoRA}\\Qwen 1.5B en 4 bits};
\node[dest] (adapter) at (10.3,0) {\textbf{Adaptador}\\pesos aprendidos};
\node[caja] (casos) at (1.9,-1.35) {\textbf{30 casos}\\solo enunciado};
\node[caja] (modelo) at (6.1,-1.35) {\textbf{Mismo Qwen}\\base / base + adaptador};
\node[caja] (salidas) at (10.3,-1.35) {\textbf{60 respuestas}\\30 por condición};
\node[dest] (juez) at (14.5,-1.35) {\textbf{FOM-5 v2}\\referencia + requisitos};
\draw[flecha] (datos)--(ajuste);
\draw[flecha] (ajuste)--(adapter);
\draw[flecha] (casos)--(modelo);
\draw[flecha] (modelo)--(salidas);
\draw[flecha] (salidas)--(juez);
\draw[flecha] (adapter.south) -- ++(0,-.21) -| (modelo.north);
\node[font=\fontsize{7}{8}\selectfont,text=azul] at (14.5,-2.1) {La referencia entra solo al evaluar};
\end{tikzpicture}
\end{center}
\vspace{-8pt}

\begin{minipage}[t]{.48\textwidth}
\small
\seccion{1}{Continuidad y elección}
La tarea de la Entrega 1 sigue intacta: traducir enunciados LP/MILP a conjuntos, parámetros, variables, objetivo y restricciones \emph{equivalentes}, sin resolver el óptimo. En el caso oftalmológico, el modelo original falló en apertura, asignación y ponderación poblacional.
\[
\begin{aligned}
\min\sum_{i,j}p_jd_{ij}x_{ij},\qquad &\sum_i y_i=2,\\[-2pt]
\sum_i x_{ij}=1\;(\forall j),\qquad &x_{ij}\le y_i\;(\forall i,j).
\end{aligned}
\]
Se eligió \textbf{Qwen2.5-1.5B-Instruct}, revisión \texttt{989aa7980e4c}, frente a Gemma-2-2B-IT (2B) y Phi-3-mini (3.8B): el menor de los tres, ejecutable en T4. Todos cometieron errores de modelación en la Entrega 1.

En una corrida por variante del caso, la rúbrica previa de localización dio 1,65/5 al prompt directo, 1,35/5 al razonamiento paso a paso y 1,00/5 a la descomposición. Se exploraron también verificación con código, selección de respuestas, one-shot y representación intermedia; persistieron fallas estructurales. Estos puntajes usan \textbf{otra rúbrica}.

\seccion{2}{Intervención y datos}
\textbf{QLoRA} entrena matrices de bajo rango sobre una base congelada y cuantizada en 4 bits (NF4). Una época usó 1050 pares sintéticos: 50 familias con 21 variantes cada una (315 LP, 735 MILP), LoRA \(r=16\), \(\alpha=32\), semilla 42 y pérdida solo sobre la respuesta. En Colab T4: \textbf{982 s, 131 pasos y 9,04 GiB} reservados.

\seccion{3}{Contraste del caso oftalmológico}
En la ejecución demostrativa, la salida pasó de \textbf{1140 a 232 tokens} (49,3 a 15,6 s). La respuesta ajustada siguió las secciones pedidas, declaró \(y_i\) y \(x_{ij}\), e incorporó cardinalidad, asignación única y vínculo \(x_{ij}\le y_i\). La respuesta base inventó una «solución óptima» y careció de \(y_i\).

\noindent\colorbox{gris}{\parbox{\dimexpr\linewidth-2\fboxsep\relax}{\small
\textbf{\color{azul}Mejora parcial.} QLoRA todavía omitió \(p_j\) del objetivo: \(\min\sum d_{ij}x_{ij}\) no minimiza distancia ponderada por población. En las respuestas archivadas de este caso, FOM-5 v2 pasó de \textbf{0,650 a 3,850/5}; el caso se reporta fuera de la media de 30.}}
\end{minipage}\hfill
\begin{minipage}[t]{.48\textwidth}
\small
\seccion{4}{Benchmark y resultados}
La comparación usa los \textbf{mismos 30 enunciados}, generación determinista y tope de 1200 tokens. Se conservaron 30 respuestas base y se generaron 30 QLoRA. Qwen3-14B NF4 juzgó ambas; referencia y requisitos entraron solo en la evaluación.

\textbf{FOM-5 v2} puntúa de 0 a 5 decisiones y dominios \(D\) (15\,\%), objetivo \(O\) (20\,\%), restricciones \(R\) (35\,\%), validez algebraica \(A\) (15\,\%) y generalización \(G\) (15\,\%):
\[
F=0{,}15D+0{,}20O+0{,}35R+0{,}15A+0{,}15G.
\]
Los requisitos ausentes o erróneos limitan las dimensiones pertinentes. La misma rúbrica evaluó ambas condiciones.

\begin{center}
\begin{tabularx}{\linewidth}{@{}Xrr@{}}
\toprule
\textbf{30 problemas} & \textbf{Base} & \textbf{QLoRA}\\
\midrule
FOM-5 v2 medio (0--5) & 0,5884 & \textbf{0,8024}\\
Casos con 0/5 & 21 & \textbf{12}\\
\bottomrule
\end{tabularx}
\end{center}
\vspace{-3pt}
\begin{center}
\begin{tikzpicture}[x=.28cm,y=.48cm,font=\fontsize{7.8}{8.5}\selectfont]
\fill[turquesa] (0,0) rectangle (14,.8);
\fill[gray!55] (14,0) rectangle (22,.8);
\fill[rojo] (22,0) rectangle (30,.8);
\node[anchor=north] at (7,-.05) {14 mejoran};
\node[anchor=north] at (18,-.05) {8 iguales};
\node[anchor=north] at (26,-.05) {8 empeoran};
\end{tikzpicture}
\end{center}
\vspace{-2pt}
El cambio medio pareado fue \textbf{+0,2140/5}. El intervalo bootstrap del 95\,\% sobre 10\,000 remuestreos de las 30 diferencias es \([{-}0,2136;\,+0,7007]\). Incluye cero: \textbf{la mejora no es concluyente} para este conjunto.

\seccion{5}{Falla y alcance}
En \texttt{opt\_02}, producción con horas extra, QLoRA bajó de \textbf{1,30 a 0,00/5}. Las horas \(h\in[0,20]\) debían ampliar solo trabajo y costar \(10h\); el modelo inventó \(y_k\), mezcló índices y omitió \(x_C\ge10\). Una plantilla aprendida sustituyó reglas del enunciado.

La concisión y el mejor formato \textbf{no garantizan corrección matemática}. Cada familia comparte una formulación canónica; tres casos tuvieron exposición histórica previa. Una semilla y un juez automático limitan la inferencia. El paso siguiente es verificar cobertura de requisitos, dominios e índices antes de aceptar un borrador.

\seccion{6}{Ejecución y materiales}
El notebook ejecuta ambas rutas sobre el mismo caso; los CSV guardan respuestas y juicios por ID. \textbf{\href{https://github.com/Jose21531/GenAI-Optimization}{Repositorio y código}} \quad \textbf{\href{https://drive.google.com/drive/folders/1NwKmSSmr4_v8Gd-W1X92eswox6rcSRxV}{Video de ejecución}}.
\end{minipage}

\vspace{3pt}
{\color{oro}\rule{\textwidth}{.8pt}}\par
{\fontsize{7.5}{8.5}\selectfont\color{azul}
Fuentes y trazabilidad: Entrega 1, rúbrica previa de localización y salidas del repositorio; Qwen2.5-1.5B-Instruct; QLoRA (Dettmers et al., 2023).}
\end{document}
